\section{Related Work}%
\label{sec:related}
According to model architectures and objective designs, self-supervised methods on graphs can be naturally divided into contrastive and generative domains. 
% Traditionally, random-walk-based network embeddings~\cite{perozzi2014deepwalk,grover2016node2vec, tang2015line} can be regarded as a simple form of generative methods with shallow neural networks to recover proximity; 
% graph kernels~\cite{shervashidze_2009_ais,shervashidze2011weisfeiler,yanardag2015deep} with designed features andmeasures can be viewed as ancestors of recent contrastive methods.

%In this section, we focus on more recent self-supervised graph learning methods based on graph neural networks (GNNs).

\subsection{Contrastive Self-Supervised Graph Learning}
Contrastive self-supervised learning
, which encourages alignment between label-invariant distributions and uniformity across other distributions, 
has been the prevalent paradigm for graph representation learning in the last two years. 
Its success relies heavily 
on the elaborate designs of the following components:

\vpara{Negative sampling.}
In pursuit of uniformity, negative sampling is a must for most contrastive methods. 
Mutual information based DGI~\cite{velivckovic2018deep} and InfoGraph~\cite{sun2019infograph} leverage corruptions to construct negative pairs.
GCC~\cite{qiu2020gcc} follows the MoCo-style~\cite{he2020momentum} negative queues. 
GRACE~\cite{zhu2020deep}, GCA~\cite{zhu2021graph} and GraphCL~\cite{you2020graph} use in-batch negatives. 
Despite recent attempts for negative-sample-free contrastive learning, strong regularization from architecture designs (e.g., BGRL ~\cite{thakoor2021bootstrapped}) or 
in-batch feature decorrelation
(e.g. CCA-SSG~\cite{zhang2021canonical}) is necessary in their practice.

\vpara{Architectures.}
Contrastive methods can be unstable in early-stage training, and thus architecture constraints are important to tackle the challenge. 
Asymmetric bi-encoder designs including momentum update~\cite{qiu2020gcc}, EMA, and Stop-gradient~\cite{thakoor2021bootstrapped} are widely adopted.

\vpara{Data augmentation.}
High-quality and informative data augmentation plays a central role in the success of contrastive learning, 
including feature-oriented (partial masking~\cite{jin2020self,hu2020strategies,you2020graph,zhu2020deep,thakoor2021bootstrapped}, shuffling~\cite{velivckovic2018deep}), proximity-oriented (diffusion~\cite{hassani2020contrastive,kefato2021self}, perturbation~\cite{you2020graph,hu2020gpt,zeng2020contrastive}), and graph-sampling-based (random-walk~\cite{qiu2020gcc,hassani2020contrastive,you2020graph}, uniform~\cite{zeng2020contrastive}, ego-network~\cite{sterling2015zinc}) augmentations. 

However, while augmentation in CV is usually human comprehensible, it is difficult to interpret in graphs. 
Without a theoretical understanding of handcrafted graph augmentation strategies, it remains unverified whether they are label-invariant and optimal. 

\subsection{Generative Self-Supervised Graph Learning}
Generative self-supervised learning aims to recover missing parts of the input data. 
It can be further classified into autoregressive and autoencoding two families. 
In previous literature, generative methods' performance on graph representation learning falls behind contrastive methods by a large margin.

\vpara{Graph autoregressive models.}
Autoregressive models decompose joint probability distributions as a product of conditionals. 
In supervised graph generation, previous researchers have proposed GraphRNN~\cite{you2018graphrnn}, GCPN~\cite{you2018graph}. 
For graph representation learning, GPT-GNN~\cite{hu2020gpt} is a recent attempt to leverage graph generation as the training objective. 
However, since most graphs do not present inherent orders, autoregressive methods make little sense on them.

\vpara{Graph autoencoders (GAEs).}
Autoencoders~\cite{hinton1994autoencoders} are designed to reconstruct certain inputs given the contexts and do not enforce any decoding orders as autoregressive methods do. 
The earliest works trace back to GAE and VGAE~\cite{kipf2016variational} which take 2-layer GCN as encoder and dot-product for link prediction decoding. EP~\cite{garcia2017learning} proposes to recover vertex features using mean squared error without input corruption. 
Later GAEs mostly adopt the structural reconstruction (e.g., ARVGA~\cite{pan2018adversarially}) following VGAE, or a combination of structural and feature reconstruction (e.g., MGAE~\cite{wang2017mgae}, GALA~\cite{park2019symmetric} and GATE~\cite{amin2020gate}) as their objectives. And NWR-GAE~\cite{tang2022graph} jointly predicts the node degree and neighbor feature distribution.

Regardless of the successful applications in link prediction and graph clustering, 
due to existing GAEs' reconstruction of structure or/and features without masking, 
their results on node/graph classification benchmarks are usually unsatisfactory.
Therefore, 
our goal in this work is 
% this work is 
to identify the weaknesses of existing GAE designs and rejuvenate the idea of self-supervised GAEs on graph representation learning for classification.

\vpara{\model v.s. Attribute-Masking.}
Recently, some works~\cite{you2020does,manessi2021graph,zhu2020self,jin2020self} have been dedicated to surveying a wide range of many self-supervision objectives' effectiveness on GNNs, including masked feature reconstruction (namely Attribute-Masking).
However, their performance lags far behind state-of-the-art contrastive methods because other critical defects of existing GAEs are not handled. 
We present a rough comparison of algorithms and results between \model and them in Appendix~\ref{app:appendix}.